\subsection{相伴覆叠}

设 B 为拓扑空间，G 为离散拓扑群，E 为 B 的以 G 为群的主覆叠。设 F 为一个 G-集；若给 F 赋予离散拓扑，则群 G 连续地作用于 F。此时 B-空间 \(E \times^{G} F\) 是一个覆叠。它称为与主覆叠 E 相伴的、以 F 为纤维型的 B 的覆叠。

定义 5. —— 设 B 为拓扑空间，G 为离散拓扑群，E 为 B 的以 G 为群的主覆叠。称 B 的覆叠 X 可伴随于主覆叠 E，如果存在一个 G-集 F 以及一个从 E × \(^{G}\) F 到 X 上的 B-同构。

设 B 为拓扑空间，G 为离散拓扑群，E 为 B 的以 G 为群的主覆叠。

对 B 的每个覆叠 X，群 G 从左边作用于 \(\mathcal{C}_{\mathrm{B}}(\mathrm{E};\mathrm{X})\) ，其法则由 \((g\cdot h)(x)=h(x\cdot g)\) 定义，其中 \(x\in E\) ，\(g\in G\) ，\(h\in\mathcal{C}_{\mathrm{B}}(\mathrm{E};\mathrm{X})\) 。

对每个赋予离散拓扑的 G-集 F，由下式定义映射 \(\alpha_{\mathrm{F}}\colon \mathrm{F}\to \mathcal{C}_{\mathrm{B}}(\mathrm{E};\mathrm{E}\times^{\mathrm{G}}\mathrm{F})\) ：

\[
\alpha_ {\mathrm{F}} (y) (x) = \pi (x, y) \quad \text {for } y \in \mathrm{F} \text { and } x \in \mathrm{E},\tag{2}
\]

其中 \(\pi\colon\mathrm{E}\times\mathrm{F}\to\mathrm{E}\times^{\mathrm{G}}\mathrm{F}\) 是典范满射。映射 \(\alpha_{\mathrm{F}}\) 与 G 在 F 上以及 \(\mathcal{C}_{\mathrm{B}}(\mathrm{E};\mathrm{E}\times^{\mathrm{G}}\mathrm{F})\) 上的作用相容。

对 B 的每个覆叠 X，存在唯一的映射 \(\beta_{X}\) （从 \(E \times^{G} \mathcal{C}_{B}(E; X)\) 到 X 中），使得：

\[
\beta_ {\mathrm{X}} (\pi (x, h)) = h (x) \quad \text {for } h \in \mathcal {C} _ {\mathrm{B}} (\mathrm{E}; \mathrm{X}) \text { and } x \in \mathrm{E}.\tag{3}
\]

事实上，有 \((g^{-1}\cdot h)(x\cdot g)=h(x)\) （对 \(x\in E\) ，\(g\in G\) 与 \(h\in\mathcal{C}_{\mathrm{B}}(\mathrm{E};\mathrm{X})\) ），这由 G 在 \(\mathcal{C}_{\mathrm{B}}(\mathrm{E};\mathrm{X})\) 上作用的定义得出。当给集合 \(\mathcal{C}_{\mathrm{B}}(\mathrm{E};\mathrm{X})\) 赋予离散拓扑时，映射 \(\beta_{X}\) 是覆叠的 B-态射。

当 \(\mathrm{X} = \mathrm{E}\times^{\mathrm{G}}\mathrm{F}\) 时，我们有

\[
\beta_ {\mathrm{X}} \left(\pi \left(x, \alpha_ {\mathrm{F}} (y)\right)\right) = \pi (x, y) \quad \text {   for   } x \in \mathrm{X} \text {   and   } y \in \mathrm{F}.\tag{4}
\]

命题 8. —— 设 B 为连通非空拓扑空间。设 G 为离散群，(E, p) 为 B 的以 G 为群的主覆叠。假设 E 连通。沿用上述记号，我们有：

\begin{enumerate}
\def\labelenumi{\alph{enumi})}
\item
  对每个赋予离散拓扑的 G-集 F，映射 \(\alpha_{F}\) 是从 G-集 F 到 G-集 \(\mathcal{C}_{\mathrm{B}}(\mathrm{E};\mathrm{E}\times^{\mathrm{G}}\mathrm{F})\) 上的同构。
\item
  设 X 为 B 的覆叠。B-态射 \(\beta_{X}\) 是从 \(E \times^{G} \mathcal{C}_{B}(E; X)\) 到 X 上的同构，当且仅当 E 的覆叠 \((E \times_{B} X, pr_{1})\) 可平凡化。
\item
  设 \(y\) 与 \(y'\) 为 F 中满足 \(\alpha_{\mathrm{F}}(y) = \alpha_{\mathrm{F}}(y')\) 的点。空间 E 非空；在其中选取一点 \(x\)。我们有 \(\pi(x, y) = \pi(x, y')\) （在 \(\mathrm{E} \times^{\mathrm{G}} \mathrm{F}\) 中）。于是存在 \(g \in \mathrm{G}\) 使得 \(x \cdot g = x\) 且 \(g^{-1} \cdot y = y'\)。前一等式蕴含 \(g\) 是 G 的单位元 \(e\) ，从而 \(y = y'\)。这样，映射 \(\alpha_{\mathrm{F}}\) 是单射。另一方面，设 \(h \in \mathcal{C}_{\mathrm{B}}(\mathrm{E}; \mathrm{E} \times^{\mathrm{G}} \mathrm{F})\) ，并设 \(x\) 为 E 的一点。设 \(x' \in \mathrm{E}\) 与 \(y' \in \mathrm{F}\) 满足 \(h(x) = \pi(x', y')\)。特别地，\(p(x) = p(x')\)；于是存在 G 的元素 \(g\) 使得 \(x' = x \cdot g\) ，并且也有 \(h(x) = \pi(x, y)\) ，其中我们置 \(y = g \cdot y'\)。B-态射 \(h\) 与 \(\alpha_{\mathrm{F}}(y)\) 在 E 的点 \(x\) 处取值相同；由于空间 E 连通，二者相等（I, 第 34 页, 命题 11 的推论 1），这就证明了映射 \(\alpha_{\mathrm{F}}\) 是满射。
\item
  由 I, 第 105 页命题 7 b) 应用于 F = \(\mathcal{C}_{\mathrm{B}}(\mathrm{E};\mathrm{X})\) ，E 上的覆叠 \(p^{*}(\mathrm{E}\times^{\mathrm{G}}\mathcal{C}_{\mathrm{B}}(\mathrm{E};\mathrm{X}))\) 同构于平凡覆叠 \(\mathrm{E}\times \mathcal{C}_{\mathrm{B}}(\mathrm{E};\mathrm{X})\) 。若 \(\beta_{\mathrm{X}}\) 是同构，则 E 上的覆叠 \(p^{*}(\mathrm{X})\) 因之可平凡化。反之，设覆叠 \(p^{*}(\mathrm{X})\) 可平凡化，我们来证明 B-态射 \(\beta_{\mathrm{X}}\) 是双射；由此将得出 \(\beta_{\mathrm{X}}\) 是 B-同构（I, 第 30 页, 命题 6 的推论 2）。
\end{enumerate}

映射 \(\beta_{\mathrm{X}}\) 是由映射 \(\gamma\colon\mathrm{E}\times\mathcal{C}_{\mathrm{B}}(\mathrm{E};\mathrm{X})\to\mathrm{X}\) （由 \(\gamma(x,h)=h(x)\) 定义）过渡到商而得的。我们来证明映射 \(\gamma\) 是满射。设 \(x\) 为 X 的一点。B-空间 E 的投影是满射，因为它是一个主覆叠；于是存在 E 的点 \(y\) 使得 \((y,x)\in\mathrm{E}\times_{\mathrm{B}}\mathrm{X}\) 。于是存在连续截面 \(s\) （可平凡化覆叠 \(\mathrm{pr}_{1}:\mathrm{E}\times_{\mathrm{B}}\mathrm{X}\to\mathrm{E}\) 的截面），使得 \(s(y)=(y,x)\) 。映射 \(h=\mathrm{pr}_{2}\circ s:\mathrm{E}\to\mathrm{X}\) 是 B-态射，且 \(h(y)=x\) ，由此得映射 \(\gamma\) 的满射性，从而得映射 \(\beta_{\mathrm{X}}\) 的满射性。

最后证明 \(\beta_{\mathrm{X}}\) 是单射。设 \((x,h)\) 与 \((x^{\prime},h^{\prime})\) 为 \(\mathrm{E}\times \mathcal{C}_{\mathrm{B}}(\mathrm{E};\mathrm{X})\) 中满足 \(h(x) = h'(x')\) 的元素。注意 \(x\) 与 \(x'\) 在 B 中有同一投影；于是存在 G 的元素 \(g\) 使得 \(x' = x\cdot g\) 。于是有 \(h(x) = h'(x\cdot g) = (g\cdot h')(x)\) 。由于空间 E 连通，有 \(h = g\cdot h'\) （I, 第 34 页, 命题 11 的推论 1）。于是 \((x,h)\) 与 \((x',h')\) 在 \(\mathrm{E}\times^{\mathrm{G}}\mathcal{C}_{\mathrm{B}}(\mathrm{E};\mathrm{X})\) 中有同一类，这表明映射 \(\beta_{\mathrm{X}}\) 是单射，证明完毕。

推论 1. —— 设 (E,p) 为 B 的以 G 为群的主覆叠；假设 E 连通且非空。B 的覆叠 X 可伴随于 E，当且仅当覆叠 \(p^{*}(X)\) 可平凡化。

若覆叠 \(p^{*}(\mathrm{X})\) 可平凡化，则由命题 8 b)，覆叠 X 可伴随于 E。此时，映射 \(\beta_{X}\) 把覆叠 X 等同于与 E 相伴的、以 \(\mathcal{C}_{\mathrm{B}}(\mathrm{E};\mathrm{X})\) 为纤维型的覆叠。反之，设 F 为赋予离散拓扑的 G-集，并设 \(X = E \times^{G} F\) 。此时 \(\alpha_{F}\) 是同构（同前引, a)），且公式 (4) 蕴含 \(\beta_{X}\) 是同构。因此覆叠 \(p^{*}(\mathrm{X})\) 可平凡化（命题 8, b)），推论得证。

推论 2. —— 设 B 为连通且局部连通的拓扑空间，(E,p) 为 B 的以 G 为群的主覆叠。设 E₀ 为 E 的一个连通分支，G₀ 为 G 中稳定 E₀ 的子群。B-空间 (E₀, p \textbar{} E₀) 是以 G₀ 为群的主覆叠（I, 第 103 页, 命题 6）。

每个可伴随于主覆叠 E 的 B 的覆叠 X 都可伴随于主覆叠 \(E_{0}\) 。特别地，覆叠 E 与主覆叠 \(E_{0}\) 相伴。

事实上，若覆叠 X 可伴随于 E，则覆叠 \(p^{*}(\mathrm{X})\) 可平凡化，从而覆叠 \(p_{0}^{*}(\mathrm{X})\) （由 \(p^{*}(\mathrm{X})\) 限制在 \(E_{0}\) 上所得）也可平凡化。

更确切地说，注意映射 \((x,g)\mapsto x\cdot g\) （从 \(E_{0}\times G\) 到 E 中）通过过渡到商，诱导出从 \(E_{0}\times^{G_{0}}G\) 到 E 上的主覆叠的 B-同构。

命题 9. —— 设 B 为拓扑空间，E 为 B 的连通非空、以 G 为群的主覆叠。设 F 为非空的、赋予离散拓扑的 G-集。为使空间 E × \(^{G}\) F 连通，充要条件是 G 传递地作用于 F。

设 U 为 \(E \times^{G} F\) 的既开又闭的子集。若以 \(\pi: E \times F \to E \times^{G} F\) 表示典范满射，则 \(\pi^{-1}(U)\) 是 \(E \times F\) 的在 G 下稳定的既开又闭子集。由于 E 连通，存在在 G 下稳定的子集 \(F' \subset F\)，使得 \(\pi^{-1}(U) = E \times F'\) ，从而 \(U = \pi(E \times F')\)。设 \(F'\) 与 \(F''\) 为 F 的在 G 下稳定的子集，满足 \(\pi(E \times F') = \pi(E \times F'')\)。由于 E 非空，有 \(F' = F''\)。映射 \(F' \mapsto \pi(E \times F')\) 是从 F 的 G 稳定子集的集合到 \(E \times^{G} F\) 的既开又闭子集的集合上的双射。命题得证。

命题 10. —— 设 B 为拓扑空间，(E,p) 为 B 的以 G 为群的主覆叠，X 为 B 的覆叠。假设 E 与 X 都连通且非空。下列性质等价：

\begin{enumerate}
\def\labelenumi{(\roman{enumi})}
\item
  覆叠 X 与主覆叠 E 相伴；
\item
  存在 G 的子群 H，使得 X B-同构于 E/H；
\item
  存在满的 B-态射 \(h: \mathrm{E} \to \mathrm{X}\) ；
\item
  对每个点 \((y, x)\) （属于 \(E \times_{B} X\) ），存在 B-态射 \(r: E \to X\) 使得 \(r(y) = x\) 。
\end{enumerate}

设这些条件成立，且 H 为使 X B-同构于 E/H 的 G 的子群。覆叠 X 是伽罗瓦的，当且仅当子群 H 在 G 中正规。

(i)⇒(ii)：设 F 为使覆叠 \(E \times^{G} F\) B-同构于 X 的离散 G-集。集合 F 非空，且空间 \(E \times^{G} F\) 连通，故群 G 传递地作用于 F（命题 9）。于是空间 \(E \times^{G} F\) B-同构于 E/H，其中 H 是 G 中稳定 F 的某个点的子群（I, 第 106 页, 例 4）。

(ii)⇒(iii)：事实上，从 E 到 E/H 上的典范满射是满的 B-态射。

(iii)⇒(iv)：设 \((y, x)\) 为 \(E \times_{B} X\) 的一点。由于映射 h 是满射，存在 E 的点 \(y'\) 使得 \(h(y') = x\)。点 y 与 \(y'\) 在 B 中有同一投影。由于覆叠 E 是以 G 为群的主覆叠，存在 \(g \in G\) 使得 \(y \cdot g = y'\)。映射 \(r: E \to X\) 由 \(z \mapsto h(z \cdot g)\) 所定义，它是 B-态射，且 \(r(y) = x\)。

(iv)⇒(i)：由于 X 非空且 E 到 B 的投影是满射，空间 \(E \times_{B} X\) 非空。于是存在从 E 到 X 的 B-态射 r。映射 ( \(Id_{E}, r\) ): \(E \to E \times_{B} X\) 是 E 的覆叠 \(p^{*}(X)\) 的一个截面。由于空间 E 连通，由 I, 第 69 页命题 1 的推论 2，覆叠 \(p^{*}(X)\) 可平凡化，这就证明覆叠 X 可伴随于伽罗瓦覆叠 p（命题 8）。

现在设这些条件成立。以 X 表示 B-空间 E/H，q 表示其投影，\(h: E \rightarrow E/H\) 表示典范映射。

若群 H 在 G 中正规，则 X 是以 G/H 为群的主覆叠（I, 第 106 页, 例 4）。由于 X 与 B 连通且非空，X 是 B 的伽罗瓦覆叠（I, 第 102 页, 定理 2）。反之，设 E/H 是 B 的伽罗瓦覆叠，并令 K = Aut \(_{B}\) (E/H) \(^{\circ}\) 。我们定义映射 \(\alpha: \mathrm{E} \times \mathrm{G} \to \mathrm{K}\) ，它把 \((t, g) \in \mathrm{E} \times \mathrm{G}\) 对应到 K 的唯一元素 \(k\) ，后者满足 \(h(t \cdot g) = h(t) \cdot k\) 。它是连续的，因为它是把连续映射 \((t, g) \mapsto (h(t), h(t \cdot g))\) （从 E \(\times\) G 到 X \(\times_{\mathrm{B}}\) X 中）与典范平凡化（I, 第 95 页, 注 1）X \(\times_{\mathrm{B}}\) X \(\to\) X \(\times\) K （即 \(q^{*}(\mathrm{X})\) 的典范平凡化）复合而得。由于 E 连通且 K 是离散群，连续映射 \(t \mapsto \alpha(t, g)\) 对每个 \(g \in \mathrm{G}\) 是常值的；我们将其值记作 \(\alpha(g)\)。若 \(t \in \mathrm{E}\) ，\(g \in \mathrm{G}\) ，且 \(g' \in \mathrm{H}\) ，则有

\[
h (t) = h \left(t \cdot g ^ {- 1} g\right) = h \left(t \cdot g ^ {- 1}\right) \cdot \alpha (g) = h \left(t \cdot g ^ {- 1} \cdot g ^ {\prime}\right) \cdot \alpha (g) = h \left(t \cdot \left(g ^ {- 1} g ^ {\prime} g\right)\right).
\]

由此得出 \(g^{-1}g'g\) 属于 H，因而 H 在 G 中正规。

推论. —— 设 E 为 B 的伽罗瓦覆叠，X 为 B 的连通非空覆叠。假设 X 是有限覆叠，或者空间 B 局部连通。若存在 B-态射 h: E → X，则覆叠 X 与主覆叠 E 相伴。

在两种情形下，B-态射 h 都是覆叠（I, 第 77 页, 定理 1 的推论 3，以及 I, 第 78 页, 命题 5），而连通空间的非空覆叠是满射（I, 第 68 页）。

定理 3. —— 设 B 为非空、连通且局部连通的拓扑空间。B 的每个覆叠都可相伴于某个伽罗瓦覆叠；B 的每个有限覆叠都可相伴于某个有限伽罗瓦覆叠。

设 X 为 B 的覆叠，以 q 表示其投影。由于空间 B 连通且非空，覆叠 X 具有次数；以 F 表示此次数，并给集合 F 赋予离散拓扑。由于空间 B 局部连通，层 \(\mathcal{I} = \underline{\mathcal{I}som}_{\mathrm{B}}(\mathrm{B} \times \mathrm{F}; \mathrm{X})\) 是局部常值的，且在 B 的每点 b 处，其茎 \(I_{b}\) 典范地同构于集合 \(\mathcal{B}(\mathrm{F}; \mathrm{X}_{b})\) ，即从 F 到纤维 \(X_{b}\) 上的双射组成的集合（I, 第 89 页, 命题 12）。设 \(E = E_{\mathcal{I}}\) 为相伴的覆叠（I, 第 86 页, 推论）。

设 G 为 F 的置换群。对每个 \(g \in G\)，我们如下定义一个态射 \(\gamma'(g)\) （层 \(\mathcal{I}\) 到自身的态射）：对 B 的每个开集 U，映射 \(\gamma'(g)_{\mathrm{U}}\) 把 U-同构 \(\varphi: \mathrm{U} \times \mathrm{F} \to q^{-1}(\mathrm{U})\) 对应到由 \((b, f) \mapsto \varphi(b, g(f))\) （其中 \(b \in U\)，\(f \in F\)）定义的 U-同构。我们有 \(\gamma'(\mathrm{Id}_{\mathrm{F}}) = \mathrm{Id}_{\mathcal{I}}\) 以及 \(\gamma'(g \circ g') = \gamma'(g') \circ \gamma'(g)\) （其中 \(g, g' \in G\)），故对每个 \(g \in G\)，\(\gamma'(g)\) 是层 \(\mathcal{I}\) 的自同构。B-态射 \(\gamma(g): E \to E\) （与 \(\gamma'(g)\) 相伴）是自同构（I, 第 50 页）。若 \(x = [\mathrm{U}, \varphi, b]\) 是 E 的元素，其中 U 是 B 的开子集，b 是 U 的点，\(\varphi\) 是从 U×F 到 \(q^{-1}(\mathrm{U})\) 上的 U-同构，则有 \(\gamma(g)(x) = [\mathrm{U}, \psi, b]\) ，其中 \(\psi\) 由 \(\psi(a, f) = \varphi(a, g(f))\) 定义（\(a \in U\)，\(f \in F\)）。若 g 与 \(g'\) 是 G 的元素，则有 \(\gamma(g \circ g') = \gamma(g') \circ \gamma(g)\)。我们有 \(\gamma(\mathrm{Id}_{\mathrm{F}}) = \mathrm{Id}_{\mathrm{E}}\)。于是，通过置 \(x \cdot g = \gamma(g)(x)\) （其中 \(x \in E\)，\(g \in G\)），我们定义了 G 在 E 上的连续右作用。群 G 单传递地作用于 E 的每条纤维，故赋予此作用的覆叠 E 是以 G 为群的主覆叠（I, 第 97 页, 命题 4）。设 h 为从 E × F 到 X 中、由 \(h([\mathrm{U}, \varphi, b], f) = \varphi(b, f)\) 定义的映射，其中 U 取遍 B 的开子集，b 取遍 U 的点，\(\varphi\) 取遍从 U × F 到 \(q^{-1}(\mathrm{U})\) 上的 U-同构。由 E 上拓扑的定义，它是连续的；它是一个 B-态射。对 G 的每个元素与 E × F 的每个点 \((x, f)\)，有 \(h(x, g(f)) = h(x \cdot g, f)\)。因此映射 h 通过过渡到商，定义一个 B-态射 \(h': E \times^{G} F \to X\)。对 B 的每个点 \(b\) ，映射 \(h_b\colon \mathrm{E}_b\times \mathrm{F}\to \mathrm{X}_b\) 等同于映射 \((\varphi, f)\mapsto \varphi(f)\) （从 \(\mathcal{B}(\mathrm{F};\mathrm{X}_b)\times \mathrm{F}\) 到 \(\mathrm{X}_b\) 中），故映射 \(h_b^\prime\) 是双射。因此，\(h^\prime\) 是从 \(\mathrm{E}\times^{\mathrm{G}}\mathrm{F}\) 到 X 上的 B-同构（I, 第 30 页, 命题 6 的推论 2）。

由于空间 B 连通、局部连通且非空，可伴随于主覆叠 E 的覆叠 X 可相伴于某个伽罗瓦覆叠（I, 第 110 页, 推论 2）。若覆叠 X 有限，则覆叠 E 亦然，故 X 可相伴于某个有限伽罗瓦覆叠（同前引）。

